\begin{table}[htbp]
\centering
\caption{Standardized Effect Sizes}
\label{tab:sde}
\begin{tabular}{lcccccc}
\hline\hline
Outcome & $\hat{\beta}$ & SE & SD($Y$) & SDE & SE(SDE) & Classification \\
\hline
Math proficiency (std.) & -0.2183 & 0.1077 & 1.00 & -0.2183 & 0.1077 & Large negative \\
\hline\hline
\end{tabular}
\begin{tablenotes}
\small
\item \textit{Notes:} \textbf{Country:} United States. \textbf{Research question:} Does predicted ground-level industrial air pollution from nearby major-emitting facilities reduce school-level math proficiency in U.S. public schools? \textbf{Policy mechanism:} Major air-emitting facilities regulated under the Clean Air Act release pollutants from industrial stacks; a Gaussian plume dispersion model predicts where emissions settle at ground level based on plant engineering parameters (stack height, meteorology) and plant-school geometry, creating within-school variation in predicted exposure driven by year-to-year wind pattern shifts that are independent of economic sorting or school investment. \textbf{Outcome definition:} School-level percentage of students scoring at or above proficiency on state math assessments (grades 3--8), from EdFacts, standardized to mean zero and unit variance within the analysis sample. \textbf{Treatment:} Continuous; sum of Gaussian plume predicted ground-level concentrations (per unit emission) from all major facilities within 50km of each school, based on Pasquill--Gifford Class D dispersion coefficients and ASOS wind direction frequencies, standardized to mean zero and unit variance. \textbf{Data:} EPA ECHO major air-emitting facility locations, Iowa Mesonet ASOS hourly wind data (2010, 2013, 2015, 2017), EdFacts school-level math proficiency (2013--2018), NCES EDGE school geocodes; school-year panel; 17,074 observations across 6,061 schools. \textbf{Method:} Reduced-form OLS with school and year fixed effects; standard errors clustered at county level (495 counties). \textbf{Sample:} U.S. public schools (grades 3--8) within 50km of at least one EPA ECHO major air-emitting facility, restricted to years with ASOS wind data. SDE $= \hat{\beta} / \text{SD}(Y)$ where SD($Y$) is the pre-treatment standard deviation. Classification refers to magnitude, not statistical significance: Large ($|$SDE$| > 0.15$), Moderate ($0.05$--$0.15$), Small ($0.005$--$0.05$), Null ($< 0.005$).
\end{tablenotes}
\end{table}
